% =====================================================================
%  Rock Size Measurement -- for the UoB `dissertation' template.
%
%  PLACEMENT: this is a \section, so % =====================================================================
%  Rock Size Measurement -- for the UoB `dissertation' template.
%
%  PLACEMENT: this is a \section, so % =====================================================================
%  Rock Size Measurement -- for the UoB `dissertation' template.
%
%  PLACEMENT: this is a \section, so % =====================================================================
%  Rock Size Measurement -- for the UoB `dissertation' template.
%
%  PLACEMENT: this is a \section, so \input{methods} it INSIDE a
%  \chapter (e.g. the Execution chapter of main.tex).
%
%  PREAMBLE: the dissertation class already loads graphicx, amsmath and
%  amssymb, so you only need to add TikZ. Put these two lines just before
%  \begin{document} in main.tex:
%
%      \usepackage{tikz}
%      \usetikzlibrary{arrows.meta, positioning, fit}
%
%  BIBLIOGRAPHY: the four \cite keys used below
%  (garrido2014aruco, yolo11, kirillov2023sam, zhang2023mobilesam)
%  are defined in references.bib -- make sure main.tex ends with
%  \bibliography{references}  (the class already sets \bibliographystyle).
% =====================================================================

\section{Rock Size Measurement}
\label{sec:measurement}

\subsection{Overview}
The size of each rock is recovered in centimetres from a single hand-held image,
using a planar reference marker and without any camera calibration. The pipeline
has four stages (Fig.~\ref{fig:pipeline}): (A) recovering the ground-plane
homography from the marker, (B) detecting the rocks, (C) segmenting their
contours, and (D) measuring each contour in metric units.

% ---------------------------------------------------------------------
%  Figure 1 -- pipeline flowchart
% ---------------------------------------------------------------------
\begin{figure}[t]
\centering
\begin{tikzpicture}[
  font=\small,
  box/.style={rectangle, draw, thin, minimum width=52mm,
              minimum height=8mm, align=center},
  term/.style={rectangle, rounded corners=3mm, draw, thin,
               minimum width=40mm, minimum height=7mm, align=center},
  arr/.style={-{Latex[length=2mm]}, thin},
  darr/.style={-{Latex[length=2mm]}, thin, dashed},
  node distance=5mm
]
  \node[term] (img)                 {Field image with marker};
  \node[box, below=of img]  (cal)   {\textbf{A}\; Detect ArUco marker;\\
                                      estimate homography $H$, scale $k$};
  \node[box, below=of cal]  (det)   {\textbf{B}\; YOLO detection;\\
                                      exclude marker, merge split boxes};
  \node[box, below=of det]  (seg)   {\textbf{C}\; SAM segmentation;\\
                                      de-duplicate masks};
  \node[box, below=of seg]  (warp)  {\textbf{D}\; Warp rock contour by $H$};
  \node[box, below=of warp] (fer)   {Maximum Feret diameter $d$};
  \node[box, below=of fer]  (size)  {$L = d/k$\quad (cm)};
  \node[box, below=of size] (cls)   {Size class (small / medium / large)};

  \draw[arr] (img)  -- (cal);
  \draw[arr] (cal)  -- (det);
  \draw[arr] (det)  -- (seg);
  \draw[arr] (seg)  -- (warp);
  \draw[arr] (warp) -- (fer);
  \draw[arr] (fer)  -- (size);
  \draw[arr] (size) -- (cls);
  % calibration supplies H, k to the measurement stage
  \draw[darr] (cal.east) to[out=0, in=0] (warp.east);
\end{tikzpicture}
\caption{Measurement pipeline. Stage~A recovers the ground-plane homography
$H$ and scale $k$ from the marker; stages~B--C localise and segment each rock;
stage~D warps the rock contour by $H$ and reads its size in centimetres. The
dashed arrow shows calibration feeding the measurement stage, either per image
or once for a fixed camera rig and reused for marker-free photos.}
\label{fig:pipeline}
\end{figure}

% ---------------------------------------------------------------------
\subsection{Ground-plane calibration}
\label{sec:calibration}
The method rests on a single geometric assumption: the reference marker and the
rocks lie on the same ground plane. Because a plane imaged by a pinhole camera is
related to any rectified view by a projective homography, the effects of camera
height, viewing angle and focal length are identical for every point on that
plane and cancel out; no intrinsic or extrinsic camera parameters are required.

A square ArUco marker~\cite{garrido2014aruco} (dictionary \texttt{DICT\_4X4\_50},
id~0) is located in the
greyscale image, with detection made robust through a three-stage cascade---raw,
CLAHE contrast-enhanced, then upsampled---and sub-pixel corner refinement. Its
four corners $p_i$ are mapped to a fronto-parallel target plane in which the
marker is a square of side $s = \mathrm{marker\_cm}\cdot k$, where the plane scale
$k = 40~\mathrm{px\,cm^{-1}}$ is fixed by construction. We solve for the
homography $H$ satisfying
\begin{equation}
  q_i \sim H\,p_i, \qquad
  q_i \in \{(0,0),\,(s,0),\,(s,s),\,(0,s)\},
  \label{eq:homography}
\end{equation}
which yields a mapping from source pixels to metric ground coordinates in which
one centimetre is exactly $k$ pixels everywhere. Two operating modes use this
result. In the \emph{per-image} mode a marker is present in every photo and $H$,
$k$ are recomputed each time. In the \emph{calibrate-once} mode a fixed camera
rig is calibrated from one marker photo; $H$ and $k$ are stored and reused for
subsequent marker-free images, valid only while the camera pose remains
unchanged.

% ---------------------------------------------------------------------
\subsection{Rock detection and box cleaning}
\label{sec:detection}
A YOLO detector~\cite{yolo11} proposes candidate rock bounding boxes. Two corrections follow.
First, the marker itself can be detected as a rock; any candidate whose
intersection-over-union with the marker bounding box exceeds $0.3$ is discarded
(skipped in the calibrate-once mode, where no marker is present). Second, the
detector occasionally splits one rock across several boxes; overlapping or
mutually contained boxes are iteratively merged into their union until the set is
stable, using an IoU threshold of $0.3$ and a containment threshold of $0.4$.
This prevents both double-counting and under-estimation of size.

% ---------------------------------------------------------------------
\subsection{Contour segmentation}
\label{sec:segmentation}
Each cleaned box is passed as a prompt to a Segment Anything model
(SAM)~\cite{kirillov2023sam}, in its lightweight MobileSAM
variant~\cite{zhang2023mobilesam}, which
returns a pixel-level mask that follows the true rock boundary far more closely
than the rectangular box. Because neighbouring prompts can produce masks of the
same rock, masks are de-duplicated: two are merged when their intersection
exceeds half the area of the smaller. When segmentation is disabled, the
rectangular detection box is used as a faster, lower-accuracy fallback.

% ---------------------------------------------------------------------
\subsection{Size measurement}
\label{sec:size}
For each rock, the largest external contour is extracted and its pixel points $c$
are warped into the metric ground plane,
\begin{equation}
  c' = H\,c .
  \label{eq:warp}
\end{equation}
The characteristic size is the maximum Feret (calliper) diameter---the largest
distance between any two points of the contour's convex hull---normalised by the
plane scale:
\begin{equation}
  L = \frac{\displaystyle\max_{i,j}\,\lVert h_i - h_j\rVert}{k}\quad[\mathrm{cm}],
  \label{eq:feret}
\end{equation}
where $\{h_i\}$ are the warped hull points. This longest-span measure is more
faithful to the true maximum diameter than the long side of a minimum-area
rectangle. Each rock is then assigned a size class from $L$: small
($<10~\mathrm{cm}$), medium ($10$--$20~\mathrm{cm}$), or large
($\geq 20~\mathrm{cm}$).

% ---------------------------------------------------------------------
\subsection{Assumptions and limitations}
\label{sec:limitations}
\begin{enumerate}
  \item \textbf{Rock-height bias.} The homography is exact only for points on the
        ground plane. A rock has thickness, so its upper surface projects to a
        displaced position and its measured size is systematically
        over-estimated. This is an intrinsic bound of single-marker monocular
        measurement and cannot be removed without depth information.
  \item \textbf{Coplanarity.} Uneven terrain, or a marker not lying flush with
        the ground, violates the single-plane assumption and degrades accuracy.
  \item \textbf{Projected metric.} $L$ is the longest diameter of the rock's
        outline projected onto the ground plane, not its true three-dimensional
        long axis.
  \item \textbf{Segmentation quality.} Under- or over-segmentation of the mask
        propagates directly into $L$.
  \item \textbf{Reference error.} An incorrect measured marker size, or
        corner-localisation error, scales proportionally into every result.
  \item \textbf{Pose drift (calibrate-once mode).} Any change in camera height,
        angle, position or zoom silently invalidates the stored $H$ and $k$.
\end{enumerate}
 it INSIDE a
%  \chapter (e.g. the Execution chapter of main.tex).
%
%  PREAMBLE: the dissertation class already loads graphicx, amsmath and
%  amssymb, so you only need to add TikZ. Put these two lines just before
%  \begin{document} in main.tex:
%
%      \usepackage{tikz}
%      \usetikzlibrary{arrows.meta, positioning, fit}
%
%  BIBLIOGRAPHY: the four \cite keys used below
%  (garrido2014aruco, yolo11, kirillov2023sam, zhang2023mobilesam)
%  are defined in references.bib -- make sure main.tex ends with
%  \bibliography{references}  (the class already sets \bibliographystyle).
% =====================================================================

\section{Rock Size Measurement}
\label{sec:measurement}

\subsection{Overview}
The size of each rock is recovered in centimetres from a single hand-held image,
using a planar reference marker and without any camera calibration. The pipeline
has four stages (Fig.~\ref{fig:pipeline}): (A) recovering the ground-plane
homography from the marker, (B) detecting the rocks, (C) segmenting their
contours, and (D) measuring each contour in metric units.

% ---------------------------------------------------------------------
%  Figure 1 -- pipeline flowchart
% ---------------------------------------------------------------------
\begin{figure}[t]
\centering
\begin{tikzpicture}[
  font=\small,
  box/.style={rectangle, draw, thin, minimum width=52mm,
              minimum height=8mm, align=center},
  term/.style={rectangle, rounded corners=3mm, draw, thin,
               minimum width=40mm, minimum height=7mm, align=center},
  arr/.style={-{Latex[length=2mm]}, thin},
  darr/.style={-{Latex[length=2mm]}, thin, dashed},
  node distance=5mm
]
  \node[term] (img)                 {Field image with marker};
  \node[box, below=of img]  (cal)   {\textbf{A}\; Detect ArUco marker;\\
                                      estimate homography $H$, scale $k$};
  \node[box, below=of cal]  (det)   {\textbf{B}\; YOLO detection;\\
                                      exclude marker, merge split boxes};
  \node[box, below=of det]  (seg)   {\textbf{C}\; SAM segmentation;\\
                                      de-duplicate masks};
  \node[box, below=of seg]  (warp)  {\textbf{D}\; Warp rock contour by $H$};
  \node[box, below=of warp] (fer)   {Maximum Feret diameter $d$};
  \node[box, below=of fer]  (size)  {$L = d/k$\quad (cm)};
  \node[box, below=of size] (cls)   {Size class (small / medium / large)};

  \draw[arr] (img)  -- (cal);
  \draw[arr] (cal)  -- (det);
  \draw[arr] (det)  -- (seg);
  \draw[arr] (seg)  -- (warp);
  \draw[arr] (warp) -- (fer);
  \draw[arr] (fer)  -- (size);
  \draw[arr] (size) -- (cls);
  % calibration supplies H, k to the measurement stage
  \draw[darr] (cal.east) to[out=0, in=0] (warp.east);
\end{tikzpicture}
\caption{Measurement pipeline. Stage~A recovers the ground-plane homography
$H$ and scale $k$ from the marker; stages~B--C localise and segment each rock;
stage~D warps the rock contour by $H$ and reads its size in centimetres. The
dashed arrow shows calibration feeding the measurement stage, either per image
or once for a fixed camera rig and reused for marker-free photos.}
\label{fig:pipeline}
\end{figure}

% ---------------------------------------------------------------------
\subsection{Ground-plane calibration}
\label{sec:calibration}
The method rests on a single geometric assumption: the reference marker and the
rocks lie on the same ground plane. Because a plane imaged by a pinhole camera is
related to any rectified view by a projective homography, the effects of camera
height, viewing angle and focal length are identical for every point on that
plane and cancel out; no intrinsic or extrinsic camera parameters are required.

A square ArUco marker~\cite{garrido2014aruco} (dictionary \texttt{DICT\_4X4\_50},
id~0) is located in the
greyscale image, with detection made robust through a three-stage cascade---raw,
CLAHE contrast-enhanced, then upsampled---and sub-pixel corner refinement. Its
four corners $p_i$ are mapped to a fronto-parallel target plane in which the
marker is a square of side $s = \mathrm{marker\_cm}\cdot k$, where the plane scale
$k = 40~\mathrm{px\,cm^{-1}}$ is fixed by construction. We solve for the
homography $H$ satisfying
\begin{equation}
  q_i \sim H\,p_i, \qquad
  q_i \in \{(0,0),\,(s,0),\,(s,s),\,(0,s)\},
  \label{eq:homography}
\end{equation}
which yields a mapping from source pixels to metric ground coordinates in which
one centimetre is exactly $k$ pixels everywhere. Two operating modes use this
result. In the \emph{per-image} mode a marker is present in every photo and $H$,
$k$ are recomputed each time. In the \emph{calibrate-once} mode a fixed camera
rig is calibrated from one marker photo; $H$ and $k$ are stored and reused for
subsequent marker-free images, valid only while the camera pose remains
unchanged.

% ---------------------------------------------------------------------
\subsection{Rock detection and box cleaning}
\label{sec:detection}
A YOLO detector~\cite{yolo11} proposes candidate rock bounding boxes. Two corrections follow.
First, the marker itself can be detected as a rock; any candidate whose
intersection-over-union with the marker bounding box exceeds $0.3$ is discarded
(skipped in the calibrate-once mode, where no marker is present). Second, the
detector occasionally splits one rock across several boxes; overlapping or
mutually contained boxes are iteratively merged into their union until the set is
stable, using an IoU threshold of $0.3$ and a containment threshold of $0.4$.
This prevents both double-counting and under-estimation of size.

% ---------------------------------------------------------------------
\subsection{Contour segmentation}
\label{sec:segmentation}
Each cleaned box is passed as a prompt to a Segment Anything model
(SAM)~\cite{kirillov2023sam}, in its lightweight MobileSAM
variant~\cite{zhang2023mobilesam}, which
returns a pixel-level mask that follows the true rock boundary far more closely
than the rectangular box. Because neighbouring prompts can produce masks of the
same rock, masks are de-duplicated: two are merged when their intersection
exceeds half the area of the smaller. When segmentation is disabled, the
rectangular detection box is used as a faster, lower-accuracy fallback.

% ---------------------------------------------------------------------
\subsection{Size measurement}
\label{sec:size}
For each rock, the largest external contour is extracted and its pixel points $c$
are warped into the metric ground plane,
\begin{equation}
  c' = H\,c .
  \label{eq:warp}
\end{equation}
The characteristic size is the maximum Feret (calliper) diameter---the largest
distance between any two points of the contour's convex hull---normalised by the
plane scale:
\begin{equation}
  L = \frac{\displaystyle\max_{i,j}\,\lVert h_i - h_j\rVert}{k}\quad[\mathrm{cm}],
  \label{eq:feret}
\end{equation}
where $\{h_i\}$ are the warped hull points. This longest-span measure is more
faithful to the true maximum diameter than the long side of a minimum-area
rectangle. Each rock is then assigned a size class from $L$: small
($<10~\mathrm{cm}$), medium ($10$--$20~\mathrm{cm}$), or large
($\geq 20~\mathrm{cm}$).

% ---------------------------------------------------------------------
\subsection{Assumptions and limitations}
\label{sec:limitations}
\begin{enumerate}
  \item \textbf{Rock-height bias.} The homography is exact only for points on the
        ground plane. A rock has thickness, so its upper surface projects to a
        displaced position and its measured size is systematically
        over-estimated. This is an intrinsic bound of single-marker monocular
        measurement and cannot be removed without depth information.
  \item \textbf{Coplanarity.} Uneven terrain, or a marker not lying flush with
        the ground, violates the single-plane assumption and degrades accuracy.
  \item \textbf{Projected metric.} $L$ is the longest diameter of the rock's
        outline projected onto the ground plane, not its true three-dimensional
        long axis.
  \item \textbf{Segmentation quality.} Under- or over-segmentation of the mask
        propagates directly into $L$.
  \item \textbf{Reference error.} An incorrect measured marker size, or
        corner-localisation error, scales proportionally into every result.
  \item \textbf{Pose drift (calibrate-once mode).} Any change in camera height,
        angle, position or zoom silently invalidates the stored $H$ and $k$.
\end{enumerate}
 it INSIDE a
%  \chapter (e.g. the Execution chapter of main.tex).
%
%  PREAMBLE: the dissertation class already loads graphicx, amsmath and
%  amssymb, so you only need to add TikZ. Put these two lines just before
%  \begin{document} in main.tex:
%
%      \usepackage{tikz}
%      \usetikzlibrary{arrows.meta, positioning, fit}
%
%  BIBLIOGRAPHY: the four \cite keys used below
%  (garrido2014aruco, yolo11, kirillov2023sam, zhang2023mobilesam)
%  are defined in references.bib -- make sure main.tex ends with
%  \bibliography{references}  (the class already sets \bibliographystyle).
% =====================================================================

\section{Rock Size Measurement}
\label{sec:measurement}

\subsection{Overview}
The size of each rock is recovered in centimetres from a single hand-held image,
using a planar reference marker and without any camera calibration. The pipeline
has four stages (Fig.~\ref{fig:pipeline}): (A) recovering the ground-plane
homography from the marker, (B) detecting the rocks, (C) segmenting their
contours, and (D) measuring each contour in metric units.

% ---------------------------------------------------------------------
%  Figure 1 -- pipeline flowchart
% ---------------------------------------------------------------------
\begin{figure}[t]
\centering
\begin{tikzpicture}[
  font=\small,
  box/.style={rectangle, draw, thin, minimum width=52mm,
              minimum height=8mm, align=center},
  term/.style={rectangle, rounded corners=3mm, draw, thin,
               minimum width=40mm, minimum height=7mm, align=center},
  arr/.style={-{Latex[length=2mm]}, thin},
  darr/.style={-{Latex[length=2mm]}, thin, dashed},
  node distance=5mm
]
  \node[term] (img)                 {Field image with marker};
  \node[box, below=of img]  (cal)   {\textbf{A}\; Detect ArUco marker;\\
                                      estimate homography $H$, scale $k$};
  \node[box, below=of cal]  (det)   {\textbf{B}\; YOLO detection;\\
                                      exclude marker, merge split boxes};
  \node[box, below=of det]  (seg)   {\textbf{C}\; SAM segmentation;\\
                                      de-duplicate masks};
  \node[box, below=of seg]  (warp)  {\textbf{D}\; Warp rock contour by $H$};
  \node[box, below=of warp] (fer)   {Maximum Feret diameter $d$};
  \node[box, below=of fer]  (size)  {$L = d/k$\quad (cm)};
  \node[box, below=of size] (cls)   {Size class (small / medium / large)};

  \draw[arr] (img)  -- (cal);
  \draw[arr] (cal)  -- (det);
  \draw[arr] (det)  -- (seg);
  \draw[arr] (seg)  -- (warp);
  \draw[arr] (warp) -- (fer);
  \draw[arr] (fer)  -- (size);
  \draw[arr] (size) -- (cls);
  % calibration supplies H, k to the measurement stage
  \draw[darr] (cal.east) to[out=0, in=0] (warp.east);
\end{tikzpicture}
\caption{Measurement pipeline. Stage~A recovers the ground-plane homography
$H$ and scale $k$ from the marker; stages~B--C localise and segment each rock;
stage~D warps the rock contour by $H$ and reads its size in centimetres. The
dashed arrow shows calibration feeding the measurement stage, either per image
or once for a fixed camera rig and reused for marker-free photos.}
\label{fig:pipeline}
\end{figure}

% ---------------------------------------------------------------------
\subsection{Ground-plane calibration}
\label{sec:calibration}
The method rests on a single geometric assumption: the reference marker and the
rocks lie on the same ground plane. Because a plane imaged by a pinhole camera is
related to any rectified view by a projective homography, the effects of camera
height, viewing angle and focal length are identical for every point on that
plane and cancel out; no intrinsic or extrinsic camera parameters are required.

A square ArUco marker~\cite{garrido2014aruco} (dictionary \texttt{DICT\_4X4\_50},
id~0) is located in the
greyscale image, with detection made robust through a three-stage cascade---raw,
CLAHE contrast-enhanced, then upsampled---and sub-pixel corner refinement. Its
four corners $p_i$ are mapped to a fronto-parallel target plane in which the
marker is a square of side $s = \mathrm{marker\_cm}\cdot k$, where the plane scale
$k = 40~\mathrm{px\,cm^{-1}}$ is fixed by construction. We solve for the
homography $H$ satisfying
\begin{equation}
  q_i \sim H\,p_i, \qquad
  q_i \in \{(0,0),\,(s,0),\,(s,s),\,(0,s)\},
  \label{eq:homography}
\end{equation}
which yields a mapping from source pixels to metric ground coordinates in which
one centimetre is exactly $k$ pixels everywhere. Two operating modes use this
result. In the \emph{per-image} mode a marker is present in every photo and $H$,
$k$ are recomputed each time. In the \emph{calibrate-once} mode a fixed camera
rig is calibrated from one marker photo; $H$ and $k$ are stored and reused for
subsequent marker-free images, valid only while the camera pose remains
unchanged.

% ---------------------------------------------------------------------
\subsection{Rock detection and box cleaning}
\label{sec:detection}
A YOLO detector~\cite{yolo11} proposes candidate rock bounding boxes. Two corrections follow.
First, the marker itself can be detected as a rock; any candidate whose
intersection-over-union with the marker bounding box exceeds $0.3$ is discarded
(skipped in the calibrate-once mode, where no marker is present). Second, the
detector occasionally splits one rock across several boxes; overlapping or
mutually contained boxes are iteratively merged into their union until the set is
stable, using an IoU threshold of $0.3$ and a containment threshold of $0.4$.
This prevents both double-counting and under-estimation of size.

% ---------------------------------------------------------------------
\subsection{Contour segmentation}
\label{sec:segmentation}
Each cleaned box is passed as a prompt to a Segment Anything model
(SAM)~\cite{kirillov2023sam}, in its lightweight MobileSAM
variant~\cite{zhang2023mobilesam}, which
returns a pixel-level mask that follows the true rock boundary far more closely
than the rectangular box. Because neighbouring prompts can produce masks of the
same rock, masks are de-duplicated: two are merged when their intersection
exceeds half the area of the smaller. When segmentation is disabled, the
rectangular detection box is used as a faster, lower-accuracy fallback.

% ---------------------------------------------------------------------
\subsection{Size measurement}
\label{sec:size}
For each rock, the largest external contour is extracted and its pixel points $c$
are warped into the metric ground plane,
\begin{equation}
  c' = H\,c .
  \label{eq:warp}
\end{equation}
The characteristic size is the maximum Feret (calliper) diameter---the largest
distance between any two points of the contour's convex hull---normalised by the
plane scale:
\begin{equation}
  L = \frac{\displaystyle\max_{i,j}\,\lVert h_i - h_j\rVert}{k}\quad[\mathrm{cm}],
  \label{eq:feret}
\end{equation}
where $\{h_i\}$ are the warped hull points. This longest-span measure is more
faithful to the true maximum diameter than the long side of a minimum-area
rectangle. Each rock is then assigned a size class from $L$: small
($<10~\mathrm{cm}$), medium ($10$--$20~\mathrm{cm}$), or large
($\geq 20~\mathrm{cm}$).

% ---------------------------------------------------------------------
\subsection{Assumptions and limitations}
\label{sec:limitations}
\begin{enumerate}
  \item \textbf{Rock-height bias.} The homography is exact only for points on the
        ground plane. A rock has thickness, so its upper surface projects to a
        displaced position and its measured size is systematically
        over-estimated. This is an intrinsic bound of single-marker monocular
        measurement and cannot be removed without depth information.
  \item \textbf{Coplanarity.} Uneven terrain, or a marker not lying flush with
        the ground, violates the single-plane assumption and degrades accuracy.
  \item \textbf{Projected metric.} $L$ is the longest diameter of the rock's
        outline projected onto the ground plane, not its true three-dimensional
        long axis.
  \item \textbf{Segmentation quality.} Under- or over-segmentation of the mask
        propagates directly into $L$.
  \item \textbf{Reference error.} An incorrect measured marker size, or
        corner-localisation error, scales proportionally into every result.
  \item \textbf{Pose drift (calibrate-once mode).} Any change in camera height,
        angle, position or zoom silently invalidates the stored $H$ and $k$.
\end{enumerate}
 it INSIDE a
%  \chapter (e.g. the Execution chapter of main.tex).
%
%  PREAMBLE: the dissertation class already loads graphicx, amsmath and
%  amssymb, so you only need to add TikZ. Put these two lines just before
%  \begin{document} in main.tex:
%
%      \usepackage{tikz}
%      \usetikzlibrary{arrows.meta, positioning, fit}
%
%  BIBLIOGRAPHY: the four \cite keys used below
%  (garrido2014aruco, yolo11, kirillov2023sam, zhang2023mobilesam)
%  are defined in references.bib -- make sure main.tex ends with
%  \bibliography{references}  (the class already sets \bibliographystyle).
% =====================================================================

\section{Rock Size Measurement}
\label{sec:measurement}

\subsection{Overview}
The size of each rock is recovered in centimetres from a single hand-held image,
using a planar reference marker and without any camera calibration. The pipeline
has four stages (Fig.~\ref{fig:pipeline}): (A) recovering the ground-plane
homography from the marker, (B) detecting the rocks, (C) segmenting their
contours, and (D) measuring each contour in metric units.

% ---------------------------------------------------------------------
%  Figure 1 -- pipeline flowchart
% ---------------------------------------------------------------------
\begin{figure}[t]
\centering
\begin{tikzpicture}[
  font=\small,
  box/.style={rectangle, draw, thin, minimum width=52mm,
              minimum height=8mm, align=center},
  term/.style={rectangle, rounded corners=3mm, draw, thin,
               minimum width=40mm, minimum height=7mm, align=center},
  arr/.style={-{Latex[length=2mm]}, thin},
  darr/.style={-{Latex[length=2mm]}, thin, dashed},
  node distance=5mm
]
  \node[term] (img)                 {Field image with marker};
  \node[box, below=of img]  (cal)   {\textbf{A}\; Detect ArUco marker;\\
                                      estimate homography $H$, scale $k$};
  \node[box, below=of cal]  (det)   {\textbf{B}\; YOLO detection;\\
                                      exclude marker, merge split boxes};
  \node[box, below=of det]  (seg)   {\textbf{C}\; SAM segmentation;\\
                                      de-duplicate masks};
  \node[box, below=of seg]  (warp)  {\textbf{D}\; Warp rock contour by $H$};
  \node[box, below=of warp] (fer)   {Maximum Feret diameter $d$};
  \node[box, below=of fer]  (size)  {$L = d/k$\quad (cm)};
  \node[box, below=of size] (cls)   {Size class (small / medium / large)};

  \draw[arr] (img)  -- (cal);
  \draw[arr] (cal)  -- (det);
  \draw[arr] (det)  -- (seg);
  \draw[arr] (seg)  -- (warp);
  \draw[arr] (warp) -- (fer);
  \draw[arr] (fer)  -- (size);
  \draw[arr] (size) -- (cls);
  % calibration supplies H, k to the measurement stage
  \draw[darr] (cal.east) to[out=0, in=0] (warp.east);
\end{tikzpicture}
\caption{Measurement pipeline. Stage~A recovers the ground-plane homography
$H$ and scale $k$ from the marker; stages~B--C localise and segment each rock;
stage~D warps the rock contour by $H$ and reads its size in centimetres. The
dashed arrow shows calibration feeding the measurement stage, either per image
or once for a fixed camera rig and reused for marker-free photos.}
\label{fig:pipeline}
\end{figure}

% ---------------------------------------------------------------------
\subsection{Ground-plane calibration}
\label{sec:calibration}
The method rests on a single geometric assumption: the reference marker and the
rocks lie on the same ground plane. Because a plane imaged by a pinhole camera is
related to any rectified view by a projective homography, the effects of camera
height, viewing angle and focal length are identical for every point on that
plane and cancel out; no intrinsic or extrinsic camera parameters are required.

A square ArUco marker~\cite{garrido2014aruco} (dictionary \texttt{DICT\_4X4\_50},
id~0) is located in the
greyscale image, with detection made robust through a three-stage cascade---raw,
CLAHE contrast-enhanced, then upsampled---and sub-pixel corner refinement. Its
four corners $p_i$ are mapped to a fronto-parallel target plane in which the
marker is a square of side $s = \mathrm{marker\_cm}\cdot k$, where the plane scale
$k = 40~\mathrm{px\,cm^{-1}}$ is fixed by construction. We solve for the
homography $H$ satisfying
\begin{equation}
  q_i \sim H\,p_i, \qquad
  q_i \in \{(0,0),\,(s,0),\,(s,s),\,(0,s)\},
  \label{eq:homography}
\end{equation}
which yields a mapping from source pixels to metric ground coordinates in which
one centimetre is exactly $k$ pixels everywhere. Two operating modes use this
result. In the \emph{per-image} mode a marker is present in every photo and $H$,
$k$ are recomputed each time. In the \emph{calibrate-once} mode a fixed camera
rig is calibrated from one marker photo; $H$ and $k$ are stored and reused for
subsequent marker-free images, valid only while the camera pose remains
unchanged.

% ---------------------------------------------------------------------
\subsection{Rock detection and box cleaning}
\label{sec:detection}
A YOLO detector~\cite{yolo11} proposes candidate rock bounding boxes. Two corrections follow.
First, the marker itself can be detected as a rock; any candidate whose
intersection-over-union with the marker bounding box exceeds $0.3$ is discarded
(skipped in the calibrate-once mode, where no marker is present). Second, the
detector occasionally splits one rock across several boxes; overlapping or
mutually contained boxes are iteratively merged into their union until the set is
stable, using an IoU threshold of $0.3$ and a containment threshold of $0.4$.
This prevents both double-counting and under-estimation of size.

% ---------------------------------------------------------------------
\subsection{Contour segmentation}
\label{sec:segmentation}
Each cleaned box is passed as a prompt to a Segment Anything model
(SAM)~\cite{kirillov2023sam}, in its lightweight MobileSAM
variant~\cite{zhang2023mobilesam}, which
returns a pixel-level mask that follows the true rock boundary far more closely
than the rectangular box. Because neighbouring prompts can produce masks of the
same rock, masks are de-duplicated: two are merged when their intersection
exceeds half the area of the smaller. When segmentation is disabled, the
rectangular detection box is used as a faster, lower-accuracy fallback.

% ---------------------------------------------------------------------
\subsection{Size measurement}
\label{sec:size}
For each rock, the largest external contour is extracted and its pixel points $c$
are warped into the metric ground plane,
\begin{equation}
  c' = H\,c .
  \label{eq:warp}
\end{equation}
The characteristic size is the maximum Feret (calliper) diameter---the largest
distance between any two points of the contour's convex hull---normalised by the
plane scale:
\begin{equation}
  L = \frac{\displaystyle\max_{i,j}\,\lVert h_i - h_j\rVert}{k}\quad[\mathrm{cm}],
  \label{eq:feret}
\end{equation}
where $\{h_i\}$ are the warped hull points. This longest-span measure is more
faithful to the true maximum diameter than the long side of a minimum-area
rectangle. Each rock is then assigned a size class from $L$: small
($<10~\mathrm{cm}$), medium ($10$--$20~\mathrm{cm}$), or large
($\geq 20~\mathrm{cm}$).

% ---------------------------------------------------------------------
\subsection{Assumptions and limitations}
\label{sec:limitations}
\begin{enumerate}
  \item \textbf{Rock-height bias.} The homography is exact only for points on the
        ground plane. A rock has thickness, so its upper surface projects to a
        displaced position and its measured size is systematically
        over-estimated. This is an intrinsic bound of single-marker monocular
        measurement and cannot be removed without depth information.
  \item \textbf{Coplanarity.} Uneven terrain, or a marker not lying flush with
        the ground, violates the single-plane assumption and degrades accuracy.
  \item \textbf{Projected metric.} $L$ is the longest diameter of the rock's
        outline projected onto the ground plane, not its true three-dimensional
        long axis.
  \item \textbf{Segmentation quality.} Under- or over-segmentation of the mask
        propagates directly into $L$.
  \item \textbf{Reference error.} An incorrect measured marker size, or
        corner-localisation error, scales proportionally into every result.
  \item \textbf{Pose drift (calibrate-once mode).} Any change in camera height,
        angle, position or zoom silently invalidates the stored $H$ and $k$.
\end{enumerate}
